\section{Principal recovery and continuation}\label{gex:recovery-section}

\subsection{The principal correction below seven eighths}

Set \(a_0=1/24000\). We extend the local estimates of
\cite[Sections 7 and 16]{OAI26} to
\begin{equation}\label{gex:second-euler-box}
 \Rea s\ge\theta,\qquad \Rea w\ge19/20,\qquad \Rea z\ge33/200.
\end{equation}
For the local ratios \(V,D,R\) of Section 7 of \cite{OAI26},
the denominators remain separated from zero:
\[
 |V|\le q_p^{-99/100},\quad |D|\le q_p^{-\theta},\quad
 |R|\le q_p^{-56/25+6a_0}.
\]
For good primes the largest defect terms have exponents
\(4-6\Rea s-6\Rea z\) and
\(1-\Rea s-\Rea w-6\Rea z\).
For row primes the six-valuation table has least decay
\(1-\Rea s-\Rea w\); the other potentially large exponents
\(3/2-3\Rea s\), \(2-4\Rea s\) are smaller. Hence
\begin{equation}\label{gex:local-euler-defect}
 |H_p-1|\ll
 \begin{cases}
 q_p^{-363/200+a_0},&p\nmid u,\\
 q_p^{-33/40+a_0},&p\mid u.
 \end{cases}
\end{equation}
The constants are uniform in phases and heights. The first product
converges normally, and the second costs \(q_u^\varepsilon\)
by the divisor-product bound. Write \(\mathcal H_{\eta,u}=\prod_{p\notin S}H_p\)
for the unselected Euler product. A fixed prime cutoff therefore gives
\begin{equation}\label{gex:principal-contraction}
 H_\eta(s)=\mathcal H_{\eta,1}(s,1,1/6)
 \text{ holomorphic on }\Rea s>\theta,\qquad
 \sup_{\Rea s>\theta}|H_\eta(s)-1|\le1/2.
\end{equation}
Further fixed exclusions preserve this bound.

The exact cancellation identity \cite[(16.8)]{OAI26} bounds
\(G_p+H_p\) by powers
\(-\theta,-99/100,-273/200+5a_0,-47/50\). At the principal
residue, where \(H_p\) is bounded away from zero, this gives
\begin{equation}\label{gex:principal-selected-factor}
 B_p=G_p/H_p=-1+O(q_p^{-4/5}).
\end{equation}
General contours still use the quotient-free tuple correction.

\subsection{Recovering the same signal}

Put
\[
 S_i(Z)=\sum_{p\in\mathcal P_i(Z)}
 W_i(q_p/P_i)q_p^{-5/6},\qquad
 A_{\mathcal T}(Z)=(-1)^K Z^{-\ell/6}\prod_iS_i(Z).
\]
The fixed-ray asymptotic of \cite[Lemma 13.1]{OAI26} gives
\[
 S_i(Z)\sim\frac{P_i^{1/6}}{|\mathcal T|\log P_i}
       \int_0^\infty W_i(y)y^{-5/6}\,dy>0.
\]
Thus, for even \(K\), \(A_{\mathcal T}(Z)>0\) eventually and
\(A_{\mathcal T}(Z)\asymp(\log Z)^{-K}\). Fixed exclusions
affect only the threshold; its inverse costs one small power.
With the unchanged positive residue constant
\[
 c_S=\widehat W_1(1)M(1/6)
       \bigl(\operatorname{Res}_{v=1}\zeta_F^S(v)\bigr)^2/6,
\]
define
\begin{equation}\label{gex:common-signal}
\begin{split}
 J_\eta(Z)&=\mathcal I_{\eta,\mathrm{mod}}(Z)/
                  \{c_S A_{\mathcal T}(Z)\},\\
 f_\eta(Z)&=\frac1{2\pi i}\int_{(2)}
 Z^{s-11/16}e^{(s-5/6)^2}
       \frac{H_\eta(s)}{L_F^S(s,\eta)}\,ds .
\end{split}
\end{equation}
Neither function contains a height cutoff.
By \eqref{gex:physical-low-bound}, for every \(\varepsilon_{\rm low}>0\),
\begin{equation}\label{gex:normalized-low}
 |J_\eta(Z)|\ll_{\eta,\varepsilon_{\rm low}}Z^{C(\theta)+\varepsilon_{\rm low}}.
\end{equation}

For \(u=1\), keep \(\Rea s=\beta_*+e\), cross \(w=1\) to
\(19/20\), then \(z=1/6\) to \(33/200\).
Repeating the residue proof of \cite[Lemma 10.5]{OAI26} in
\eqref{gex:second-euler-box} crosses only the two scalar zeta poles.
Its geometric savings are
\[
 m_w=l_y/20=5749/240000,\qquad
 m_z=h/600=3251/2400000.
\]
The latter includes all selected-slot powers.
At the double residue,
\(X^{1/3}Z^{s-5/6}=Z^{-\ell/6}Z^{s-11/16}\);
\eqref{gex:principal-selected-factor} gives relative error
\(O(Z^{-\kappa_P})\), with
\(\kappa_P=(2/5)\min_i\ell_i>0\).
Move the main integral right to \(2\), bounding the slot error
on its original line. The Gaussian pays the horizontal joins.
The principal recovery error is
\begin{equation}\label{gex:principal-recovery-error}
 \ll Z^{C(\beta_*)+(1+h)e-m_w+\varepsilon_{\rm pr}}
 +Z^{C(\beta_*)+e-m_z+\varepsilon_{\rm pr}}
 +Z^{C(\beta_*)+e-\kappa_P+\varepsilon_{\rm pr}}.
\end{equation}

\subsection{Outer rows and height closure}

For \(d<1/100\), use \(\Rea s=7/8+8e\), to the right of
\(\beta_*\), and \((\Rea w,\Rea z)=(1/2,17/50)\).
The full absolute tuple bound, including selected primes dividing
\(u\), is \(q_u^\varepsilon Z^{17\ell/50}\); the row exponent
is \(63/50<2\). Paying \(7/8-\beta_*<p/4\) leaves saving
\[
 m_{\rm small}=63/800-51p/100-p/4>1/16.
\]
This includes bounded nontrivial units and principal denominators.
In the absolute estimate, \(G_p=O(1)\) off \(u\) and
\(G_p=O(q_p^{1/2})\) on \(u\). The latter choices are
divisor-bounded, and their actual weight
\(q_p^{\Rea z-1}\) absorbs this growth
\cite[Lemma 16.2]{OAI26}. These ramified selections therefore
remain in the tuple sum.
For the full large-row contributions \(\mathscr R_\eta(U;Z)\)
in \eqref{gex:high-identity}, absolute lines \((2,2,z_\infty)\) give
\[
 \sum_{U>Z^{h+\zeta}}|\mathscr R_\eta(U;Z)|
 \ll Z^{B_0+(h+\zeta)(1+\varepsilon)-\zeta z_\infty+\varepsilon},
 \qquad B_0=l_x/2+1+l_y.
\]
Choose \(z_\infty\) after \(\zeta\) to obtain the required saving.

Choose \(\varepsilon_{\rm low}=\Gamma/2\). Before fixing a target, choose moment
losses, the mesh and disjoint slots, then the amplitude and contour
widths so that \eqref{gex:central-loss-budget} holds and
\[
 (1+h)e+\varepsilon_{\rm pr}<m_w/4,\qquad
 e+\varepsilon_{\rm pr}<\min\{m_z,\kappa_P\}/4.
\]
Reserve fractions of the remaining savings for normalizing and
summing dyads. Choose \(\zeta,z_\infty\) and the prime cutoff next.
For each target fix its arithmetic data and internal seminorm and
Sobolev orders. The common height allocation of
\cite[Sections 10 and 20.6]{OAI26} then gives
\[
 |J_\eta-f_\eta|\ll_{\eta,N}
 Z^{C(\beta_*)-m_1}(1+T_1)^{A_\eta}+Z^{B_\eta}T_1^{-N},
 \qquad T_1=Z^{\tau_\eta} ,
\]
where \(m_1>0\) is common to all targets, and \(A_\eta,B_\eta\)
are finite and independent of the later external order \(N\).
Choose \(\tau_\eta>0\) to meet the detector ceiling and
\(\tau_\eta A_\eta<m_1/4\), then \(N_\eta\) large enough to pay
all discarded axes and joins, and finally the threshold.
Increasing external orders changes only constants and thresholds.
Thus some common \(\sigma>0\) satisfies
\begin{equation}\label{gex:common-low-high}
 |J_\eta|\ll_\eta Z^{C(\theta)+\varepsilon_{\rm low}},\qquad
 |J_\eta-f_\eta|\ll_\eta Z^{C(\beta_*)-\sigma}.
\end{equation}

\subsection{Continuation}

Set \(\epsilon_*=\min\{\Gamma-\varepsilon_{\rm low},\sigma\}>0\).
Then \(f_\eta(Z)=O_\eta(Z^{C(\beta_*)-\epsilon_*})\) at infinity.
At zero it decays faster than every fixed power, by shifting its
integral to any \(B>2\), using the bounded reciprocal Euler product,
\eqref{gex:principal-contraction}, and the Gaussian.
Consequently
\[
 F_\eta(s)=\int_0^\infty f_\eta(Z)Z^{-C(s)}\,\frac{dZ}{Z}
\]
is holomorphic for \(\Rea s>\beta_*-\epsilon_*>\theta\).
Fourier inversion with \(Z=e^u\) identifies it on \(\Rea s>1\)
as \(e^{(s-5/6)^2}H_\eta(s)/L_F^S(s,\eta)\).
Since \(|H_\eta|\ge1/2\), this continues the reciprocal
holomorphically to the common half-plane
\(\Rea s>\beta_*-\epsilon_*\).
Deleted Euler factors are nonzero there, while the definition of
\(\beta_*\) supplies a primitive target with a zero in that
half-plane. This contradiction proves \(\beta_*\le\theta\).
The saving \(\epsilon_*\) is common to the family; it is precisely
this uniformity, rather than a common lower threshold or implied
constant, that lets us use the zero supremum. The argument places
no bound on the imaginary part of the excluded zero.

This proves the Hecke assertion for primitive characters.
An induced presentation differs by finitely many factors
\(1-\eta_0(\mathfrak p)q_{\mathfrak p}^{-s}\), which are nonzero
for \(\Rea s>0\). The same boundary therefore holds for all
finite-order Hecke presentations.

For a Dirichlet character \(\chi\), form the finite-order Hecke
character \(\eta_\chi(\mathfrak a)=\chi(\N\mathfrak a)\).
Delete a finite rational prime set \(S_{\Q}\) containing \(3\)
and the modulus support, and let \(S_F\) consist of all primes
above it. The split and inert Euler factors give
\[
 L_F^{S_F}(s,\eta_\chi)
 =L^{S_{\Q}}(s,\chi)L^{S_{\Q}}(s,\chi\chi_{-3}).
\]
On \(\theta<\Rea s<1\), both factors on the right are
holomorphic and their product is nonzero. Restoring the finite
Euler factors proves Dirichlet nonvanishing there.
For \(\Rea s\ge1\), use the existing theorem of \cite{OAI26},
with principal poles excepted. This also covers imprimitive
Dirichlet characters and, in particular, the zeta function.

Finally, the functional equations reflect a critical-strip zero
of a primitive character to a zero of its conjugate character
with real part \(1-\Rea\rho\). The family is closed under
conjugation. The deleted factors have no zeros in the critical
strip, so every member has the lower boundary
\(1-\theta=3001/24000\). This proves Theorem~\ref{gex:main}.
